% Recovery of 13_accion_interaccion_absorbedor and 83_operador_absorbedor.
% Development of the chronological channel from route transport.
% Provenance and conditions in PROCEDENCIA_PROPAGACION_BIDIRECCIONAL.md.
% The originals and section 80 remain intact. No radion content.

\section{Bidirectional propagation and the absorption condition at the boundary}
\label{vii:sec:propagacion-bidireccional}

APP--TRIT--TPK transport preserves the orientation of the route and the memory
that distinguishes its outward passage from its return. On this antecedent,
three successive operations are constructed: the action of an enriched history,
the propagation of a source through the realized transports, and the imposition
of a common condition at the boundary. The inscription of the record in an
apparatus space, established in Section~\ref{vii:sec:observador-aparato}, allows
these operations to be expressed through linear maps without conflating the
observable representation with the state that preserves its antecedents.

Temporal direction, sheet conjugation, and spatial reflection have different
domains. Their relations are first specified on the record; the intertwinings
used by a dynamical realization are then specified. This order determines what
information remains in the propagators and what condition is introduced by the
choice of an absorbing boundary.

\subsection{Route action and involutions of the record}
\label{vii:subsec:accion-bidireccional}

A finite route contains its positions, cardinal directions
\((d_1,\ldots,d_N)\), sheets \((\mu_1,\ldots,\mu_N)\), with
\(\mu_t\in\{\Sigma,\Pi\}\), and the increments of its enriched
record. For \(\lambda\geq0\), we define
\begin{equation}
 \mathcal A_\lambda(\gamma)
 =\sum_{t=2}^N
 \bigl(1-\delta_{d_t,d_{t-1}}
       +\lambda(1-\delta_{\mu_t,\mu_{t-1}})\bigr).
 \label{vii:eq:accion-ruta-bidireccional}
\end{equation}
The first summand counts changes of direction; the second counts changes of
sheet with the stated weight. For \(N\leq1\), the sum is zero.
The general action of the sector also retains its boundary and memory terms:
\begin{equation}
 \mathcal S[\gamma]
 =\sum_{a\subset\gamma}\mathcal L_a
  +\mathcal B_{\rm in}(\partial_-\gamma)
  +\mathcal B_{\rm out}(\partial_+\gamma)
  +\mathcal M(\gamma).
 \label{vii:eq:accion-enriquecida-frontera}
\end{equation}
The terms in this expression specify a variational class with typed endpoints
and preserved memory. The dynamics of the sector determines
\(\mathcal L_a\), the admissible conditions, and the term \(\mathcal M\).

On the records we define \(\mathcal C_{\rm r}\) as the global exchange of
sheets; \(\mathcal P_{\rm r}\) as a reflection of the spatial chart and its
directions; and \(\mathcal T_{\rm r}\) as the reversal of temporal order,
the replacement of each direction by its opposite, and the exchange of
endpoints. Cardinal reflection is an involutive permutation.
Reversal of the signed record is
\(C_M(a_1,\ldots,a_N)=(-a_N,\ldots,-a_1)\).
The sheets, quotients, and memory are transported with the complete route;
the scalar reading of a word is a subsequent operation.

\begin{proposition}[Invariance of the local route action]
\label{vii:prop:invariancia-ruta-bidireccional}
On the preceding domain of records,
\[
 \mathcal A_\lambda(\gamma)
 =\mathcal A_\lambda(\mathcal C_{\rm r}\gamma)
 =\mathcal A_\lambda(\mathcal P_{\rm r}\gamma)
 =\mathcal A_\lambda(\mathcal T_{\rm r}\gamma)
 =\mathcal A_\lambda(\mathcal C_{\rm r}\mathcal P_{\rm r}
                              \mathcal T_{\rm r}\gamma).
\]
\end{proposition}
\begin{proof}
A global permutation of the alphabet of directions preserves equality
between consecutive directions. Exchanging the two sheets preserves equality
between consecutive sheets. Reversing the order establishes a bijection between
adjacent pairs, and replacing directions by their opposites also preserves their
equalities. Each operation preserves the two counts in
\eqref{vii:eq:accion-ruta-bidireccional}; their composition preserves them.
\end{proof}

The proposition determines the symmetry of the local functional. For a
transformation also to act on a dynamical problem with fixed boundaries,
it must transport its class of admissible prolongations and the other three
terms in \eqref{vii:eq:accion-enriquecida-frontera}. In particular, two
routes with the same final representation may retain different variational
classes because of their memory.

An interaction between two sections realized at a boundary \(B\)
uses a coupling map
\[
 \mathcal I_B:E_B^{(1)}\otimes E_B^{(2)}\longrightarrow E_B^{(12)}
\]
and a term \(\mathcal S_{\rm int}\) in the same action. Its covariance
is expressed by
\[
 \rho_{12}(g)\mathcal I_B
 =\mathcal I_B\bigl(\rho_1(g)\otimes\rho_2(g)\bigr).
\]
When the histories are composed, the record of the shared boundary is counted
once; the memories of both sides are transported to the composite section.
The coupling and this covariance condition are data of the specific interaction,
distinct from the involutions of the record.

\subsection{Reversal of the realized evolution}
\label{vii:subsec:reversion-evolucion}

The realization of the record provides a real Hilbert space
\(\mathcal H_{\mathbb R}\), an orthogonal complex structure
\(J_E\), with \(J_E^2=-I\), and a self-adjoint Hamiltonian \(H\).
Its domain \(\mathcal D(H)\) is invariant under \(J_E\); we assume
\(HJ_E=J_EH\) on that domain. Let \(\mathcal C\) be an orthogonal
involution that preserves \(\mathcal D(H)\) and satisfies
\begin{equation}
 \mathcal C J_E\mathcal C^{-1}=-J_E,\qquad
 \mathcal C H\mathcal C^{-1}=H.
 \label{vii:eq:dominio-reversion-evolucion}
\end{equation}
The internal action constant \(\hbar_{\rm int}>0\) retains its
previous generation; here it normalizes the evolution group.

\begin{theorem}[Temporal conjugation of the group]
\label{vii:thm:reversion-evolucion}
The generator \(A=-J_EH/\hbar_{\rm int}\), with domain \(\mathcal D(H)\),
is skew-adjoint. Its orthogonal group, unitary with respect to \(J_E\), satisfies
\begin{equation}
 U(t)=\exp(-J_EHt/\hbar_{\rm int}),\qquad
 \mathcal C U(t)\mathcal C^{-1}=U(-t).
 \label{vii:eq:reversion-grupo-bidireccional}
\end{equation}
\end{theorem}
\begin{proof}
Orthogonality and \(J_E^2=-I\) imply \(J_E^*=-J_E\).
Domain invariance and commutation with \(H\) allow
\(H\) to be treated as a complex self-adjoint operator in the space defined by
\(J_E\). By functional calculus, \(-J_EH/\hbar_{\rm int}\) is
skew-adjoint and generates the stated group. In finite dimension, the same
assertions follow from the exponential series and \(A^*=-A\).
The hypotheses on \(\mathcal C\) give
\(\mathcal C A\mathcal C^{-1}=-A\). Conjugating the group, or its
functional calculus, yields \(\exp(-tA)=U(-t)\).
\end{proof}

The map realizing the signed involution \(C_M\) as \(\mathcal C\)
preserves temporal orientation and the preceding domain. Charge conjugation
in a field representation also preserves its own characters and operators.
Equality \eqref{vii:eq:reversion-grupo-bidireccional} specifies the temporal
reversal of the group under the stated conditions.

\subsection{Propagators of the chronological channel of a route}
\label{vii:subsec:green-cronologico}

The isometric inscription of the record admits the unitary prolongations
of Theorem~\ref{vii:thm:observador-entrelazador}. Fix a bilateral realization
of a chronological channel with fibres \(\mathcal H_n\),
\(n\in\mathbb Z\), and unitary transports
\(U_n:\mathcal H_n\to\mathcal H_{n+1}\).
The choice includes auxiliary spaces when required; it preserves
the route and its preceding records. We denote by \(U(n,k)\) the transport
from \(k\) to \(n\), with
\[
 U(k,k)=I,\qquad U(n,k)U(k,l)=U(n,l),\qquad
 U(n,k)^*=U(k,n).
\]
For \(n>k\), it is the product \(U_{n-1}\cdots U_k\); for \(n<k\),
it is the inverse of the transport from \(n\) to \(k\).

The source space \(\mathcal J_c\) contains the finitely supported sections.
The response space \(\mathcal F\) contains all sections of these fibres.
The quadratic route action, in the internal units of this realization, uses
\[
 (D\psi)_n=\psi_{n+1}-U_n\psi_n,\qquad
 E[\psi]=\sum_n\|(D\psi)_n\|^2
\]
for finitely supported sections. The kinetic operator with the sign convention
\(L=-D^*D\) is
\begin{equation}
 (L\psi)_n=U_n^*\psi_{n+1}-2\psi_n+U_{n-1}\psi_{n-1}.
 \label{vii:eq:diferencia-covariante-bidireccional}
\end{equation}
Indeed, \((D^*z)_n=z_{n-1}-U_n^*z_n\), and substitution yields
this formula. The variation of \(E\) is
\(-2\operatorname{Re}\langle\delta\psi,L\psi\rangle\).

\begin{theorem}[Exact Green operators of the covariant difference]
\label{vii:thm:green-cronologico}
Writing \(x_+=\max(x,0)\), the maps
\begin{align}
 (\mathsf{G}_{\rm ret}j)_n&=\sum_k(n-k)_+U(n,k)j_k,\label{vii:eq:green-ret-cronologico}\\
 (\mathsf{G}_{\rm adv}j)_n&=\sum_k(k-n)_+U(n,k)j_k\label{vii:eq:green-adv-cronologico}
\end{align}
from \(\mathcal J_c\) to \(\mathcal F\) satisfy
\[
 L\mathsf{G}_{\rm ret}j=L\mathsf{G}_{\rm adv}j=j.
\]
The retarded response vanishes before and at the first source position;
the advanced response vanishes after and at the last. They are the unique
solutions that vanish sufficiently far in the respective direction.
For finitely supported \(\psi\), one also has
\(\mathsf{G}_{\rm ret}L\psi=\mathsf{G}_{\rm adv}L\psi=\psi\).
With respect to pairing with finitely supported sources,
\(\mathsf{G}_{\rm adv}=\mathsf{G}_{\rm ret}^{\dagger}\).
\end{theorem}
\begin{proof}
Each sum has as many terms as the source has positions.
Composition of the transports reduces application of \(L\) to the scalar
identity
\[
 (n+1-k)_+-2(n-k)_++(n-1-k)_+=\delta_{nk}.
\]
The function \((k-n)_+\) satisfies the same identity. This proves
both equations and the support statements. For uniqueness, a homogeneous
solution that vanishes at two consecutive positions is determined to be zero
at the next by \eqref{vii:eq:diferencia-covariante-bidireccional};
the recurrence, using invertible transports, extends in both
directions. The difference of two solutions with the same causal support
vanishes sufficiently far away and is therefore zero.
Applying this argument to \(\psi\) and \(\mathsf{G}_{\rm ret}L\psi\), or to the
advanced case, gives the right-hand identities on finite support.
Finally, for two compactly supported sources, the double sum in the pairing
is finite and
\[
 \bigl((k-n)_+U(k,n)\bigr)^*=(k-n)_+U(n,k).
\]
Interchanging the indices proves the stated adjoint relation.
\end{proof}

The responses belong to \(\mathcal F\); in general they grow linearly
outside the source and do not belong to \(\ell^2\).
The finite-support domain and its pairing specify the Green adjoint relation
used here. The result constructs propagation in the chronological channel
from its action; a spatial-field realization additionally retains its own
operator, geometry, and propagation domain.

Let \(C_0\) be an isometric involution, linear or antilinear, of
\(\mathcal H_0\). Transport from the reference fibre constructs
\[
 C_n=U(-n,0)C_0U(0,n):\mathcal H_n\longrightarrow\mathcal H_{-n},
 \qquad (\mathcal C\psi)_n=C_{-n}\psi_{-n}.
\]
By composition, \(C_{-n}C_n=I\) and
\(C_{n+1}U_n=U_{-n-1}^*C_n\).
In particular, \(\mathcal C^2=I\), and the substitution
\((n,k)\mapsto(-n,-k)\) in the preceding sums proves
\begin{equation}
 \mathcal C L=L\mathcal C,\qquad
 \mathcal C \mathsf{G}_{\rm ret}\mathcal C^{-1}=\mathsf{G}_{\rm adv}.
 \label{vii:eq:intercambio-green-cronologico}
\end{equation}
When \(C_0\) is the stated realization of the record involution,
these formulas transport that involution to the entire channel, without
identifying phase return with a memory reset.

\subsection{Causal interchange and decomposition of the response}
\label{vii:subsec:green-realizacion-campos}

In a field realization, let
\(L:\mathcal D(L)\subset\mathcal H_\Phi\to\mathcal H_J\)
be an operator with unique retarded and advanced Green operators on the
admissible source domains. We denote by \(J^+(S)\) and \(J^-(S)\) the
future and past causal domains of a support \(S\). The conditions are
\[
 \begin{gathered}
 L\mathsf{G}_{\rm ret}=L\mathsf{G}_{\rm adv}=I,\\
 \operatorname{supp}(\mathsf{G}_{\rm ret}j)\subset J^+(\operatorname{supp}j),\\
 \operatorname{supp}(\mathsf{G}_{\rm adv}j)\subset J^-(\operatorname{supp}j).
 \end{gathered}
\]
Let \(C_\Phi,C_J\) be isometric involutions, both linear or both
antilinear, that transport these domains, reverse their causal orientations,
and satisfy \(LC_\Phi=C_JL\). A Green pairing is retained for which
\(\mathsf{G}_{\rm adv}=\mathsf{G}_{\rm ret}^{\dagger}\).

\begin{proposition}[Conjugation of the Green operators and parities]
\label{vii:prop:green-paridades}
Under the preceding conditions,
\[
 C_\Phi \mathsf{G}_{\rm ret}C_J^{-1}=\mathsf{G}_{\rm adv}.
\]
On the real space of response operators, define
\[
 \Theta(K)=C_\Phi K C_J^{-1},\qquad
 \mathsf P_\pm(K)=\tfrac12(K\pm\Theta(K)).
\]
These are complementary projections for the involution \(\Theta\), and
\begin{equation}
 \mathsf{G}_{\rm sym}=\tfrac12(\mathsf{G}_{\rm ret}+\mathsf{G}_{\rm adv}),\qquad
 \mathsf{G}_{\rm rad}=\tfrac12(\mathsf{G}_{\rm ret}-\mathsf{G}_{\rm adv})
 \label{vii:eq:descomposicion-green}
\end{equation}
satisfy \(\Theta(\mathsf{G}_{\rm sym})=\mathsf{G}_{\rm sym}\),
\(\Theta(\mathsf{G}_{\rm rad})=-\mathsf{G}_{\rm rad}\),
\(\mathsf{G}_{\rm sym}^{\dagger}=\mathsf{G}_{\rm sym}\), and
\(L\mathsf{G}_{\rm rad}=0\). The support of the symmetric response is contained
in the union of the two causal domains.
\end{proposition}
\begin{proof}
The conjugated operator solves the source equation, since
\[
 LC_\Phi \mathsf{G}_{\rm ret}C_J^{-1}=C_JL\mathsf{G}_{\rm ret}C_J^{-1}=I.
\]
Reversal of orientation carries its support to the advanced domain;
uniqueness identifies it with \(\mathsf{G}_{\rm adv}\). The involutions give
\(\Theta^2=I\); expanding \((I\pm\Theta)^2\) and
\((I+\Theta)(I-\Theta)\) proves the projection assertions.
The Green adjoint relation proves symmetry, and subtracting the two source
equations gives \(L\mathsf{G}_{\rm rad}=0\). The sum of two responses has support
contained in the union of their supports.
\end{proof}

If the involutions are antilinear, \(\Theta\) is linear over real
scalars; this is the operator space used in the proposition.
Conjugation of the temporal group, the support condition, and the Green
adjoint relation are three distinct properties, transmitted by the
corresponding realization maps.

\subsection{Projection of absorbed currents and variation of the action}
\label{vii:subsec:proyeccion-absorcion}

Let \(\gamma_\partial:\mathcal H_\Phi\to\mathcal H_\partial\)
be the boundary reading, distinct from the route \(\gamma\). Define
\[
 \mathcal A_\partial=\gamma_\partial \mathsf{G}_{\rm rad}.
\]
The finite window fixes the source space and boundary reading;
the Green operators retain their response domain. At a finite level with
fixed inner products, its Moore--Penrose inverse determines
\begin{equation}
 P_{\rm abs}=I-\mathcal A_\partial^+\mathcal A_\partial.
 \label{vii:eq:proyector-absorcion}
\end{equation}

\begin{theorem}[Absorption condition and boundary response]
\label{vii:thm:proyeccion-absorcion}
\(P_{\rm abs}\) is the orthogonal projection onto
\(\ker\mathcal A_\partial\). For any source \(j\), the source
\(j_{\rm abs}=P_{\rm abs}j\) satisfies
\begin{equation}
 \gamma_\partial \mathsf{G}_{\rm ret}j_{\rm abs}
 =\gamma_\partial \mathsf{G}_{\rm adv}j_{\rm abs}.
 \label{vii:eq:igualdad-borde-absorbido}
\end{equation}
It is the absorbed source closest to \(j\) for the stated inner product.
Nonzero absorbed sources exist exactly when
\(\operatorname{rank}\mathcal A_\partial<\dim\mathcal H_J\).
\end{theorem}
\begin{proof}
The orthogonal decomposition
\(\mathcal H_J=\ker\mathcal A_\partial\oplus
 \operatorname{im}\mathcal A_\partial^*\) shows that
\(\mathcal A_\partial^+\mathcal A_\partial\) is the identity on the
second summand and zero on the first. Its complement is the stated
projection. By definition,
\(2\mathcal A_\partial j_{\rm abs}
 =\gamma_\partial(\mathsf{G}_{\rm ret}-\mathsf{G}_{\rm adv})j_{\rm abs}=0\).
For \(k\in\ker\mathcal A_\partial\), the orthogonal decomposition gives
\[
 \|j-k\|^2=\|(I-P_{\rm abs})j\|^2+\|P_{\rm abs}j-k\|^2.
\]
The unique minimum corresponds to \(k=P_{\rm abs}j\).
The final equivalence is the rank--nullity theorem.
\end{proof}

In the source--field pairing, the symmetric action is
\(\mathscr S[j]=\tfrac12\langle j,\mathsf{G}_{\rm sym}j\rangle\).
The stated adjoint relation makes this quantity real. For real variations,
\[
 \mathrm d\mathscr S[j](h)
 =\operatorname{Re}\langle h,\mathsf{G}_{\rm sym}j\rangle.
\]
The action restricted to absorbed sources is written on the ambient space as
\begin{equation}
 \mathscr S_{\rm abs}[j]
 =\tfrac12\langle P_{\rm abs}j,\mathsf{G}_{\rm sym}P_{\rm abs}j\rangle,
 \quad
 \mathrm d\mathscr S_{\rm abs}[j](h)
 =\operatorname{Re}\langle P_{\rm abs}h,\mathsf{G}_{\rm sym}P_{\rm abs}j\rangle.
 \label{vii:eq:variacion-restringida-absorcion}
\end{equation}
When the pairing identifies field and source through Riesz representation,
the ambient gradient is \(P_{\rm abs}\mathsf{G}_{\rm sym}P_{\rm abs}j\).
The field response remains \(\mathsf{G}_{\rm sym}j_{\rm abs}\); the gradient
projection expresses the restriction of variations.
The formulas are proved by expanding the quadratic term of
\(j+th\) and using self-adjointness. A matter term
\(\mathscr S_{\rm matter}\) adds its own differential.

\subsection{Explicit absorbing condition in a chronological window}
\label{vii:subsec:absorcion-momentos}

The channel in Subsection~\ref{vii:subsec:green-cronologico} allows
this condition to be evaluated without introducing a target response.
Let a source be supported on \(0,\ldots,N\), with \(N\geq1\).
Transport its components to the reference fibre:
\[
 \widehat j_k=U(0,k)j_k,\qquad
 M_0(j)=\sum_{k=0}^N\widehat j_k,\qquad
 M_1(j)=\sum_{k=0}^Nk\widehat j_k.
\]
The identity \((n-k)_+-(k-n)_+=n-k\) gives
\begin{equation}
 (\mathsf{G}_{\rm ret}-\mathsf{G}_{\rm adv})j\big|_n
 =U(n,0)\bigl(nM_0(j)-M_1(j)\bigr).
 \label{vii:eq:diferencia-green-momentos}
\end{equation}

\begin{proposition}[Absorption and two transported moments]
\label{vii:prop:absorcion-momentos}
If the trace evaluates the field at two distinct positions \(a,b\),
equality of the advanced and retarded traces is equivalent to
\(M_0(j)=M_1(j)=0\). In a fibre of dimension \(d\), the space of
such sources has dimension \((N-1)d\). For \(N\geq2\), it
contains the nonzero sources with transported components
\((z,-2z,z,0,\ldots,0)\), \(z\neq0\).
\end{proposition}
\begin{proof}
Invertibility of \(U(a,0)\) and \(U(b,0)\) reduces the boundary
condition to \(aM_0-M_1=0\) and \(bM_0-M_1=0\).
Subtracting the equations gives \((a-b)M_0=0\), and then \(M_1=0\).
The two-moment map is surjective: for prescribed values
\((z_0,z_1)\), it suffices to take
\(\widehat j_0=z_0-z_1\), \(\widehat j_1=z_1\), and all other
components zero. Its rank is \(2d\); its nullity is
\((N+1)d-2d\). The indicated source has both moments zero.
\end{proof}

In the transported frame, write
\[
 B_N=\begin{pmatrix}1&1&\cdots&1\\0&1&\cdots&N\end{pmatrix}.
\]
The absorbing projection is evaluated by
\begin{equation}
 P_N=\bigl[I-B_N^{\mathsf T}(B_NB_N^{\mathsf T})^{-1}B_N\bigr]
       \otimes I_{\mathcal H_0}.
 \label{vii:eq:proyector-momentos-cronologicos}
\end{equation}
The two rows of \(B_N\) are independent for \(N\geq1\), so the
inverse exists. Unitary transport between frames carries this projection
to the original fibres. For \(N=2\),
\[
 P_2=\frac16
 \begin{pmatrix}1&-2&1\\-2&4&-2\\1&-2&1\end{pmatrix}
 \otimes I_{\mathcal H_0}.
\]
This construction produces a nonzero absorbed class and its projection
operation in the chronological channel. Its amplitude and its identification
with a matter current retain the rules of the corresponding physical sector.

\subsection{Refinement and continuity of the boundary condition}
\label{vii:subsec:refinamiento-absorcion}

Let \(R_m^\Phi,R_m^J,R_m^\partial\) be the refinement maps for
fields, sources, and boundaries. Suppose that they intertwine both Green
operators and the trace. Composition gives
\begin{equation}
 R_m^\Phi \mathsf{G}_{{\rm sym},m}=\mathsf{G}_{{\rm sym},m+1}R_m^J,\qquad
 R_m^\partial\mathcal A_{\partial,m}
 =\mathcal A_{\partial,m+1}R_m^J.
 \label{vii:eq:refinamiento-lectura-absorcion}
\end{equation}
The second equality carries an absorbed source to another absorbed source.
Naturality of the projections additionally preserves the orthogonal structure
of the realizations.

\begin{proposition}[Naturality of the absorbing projection]
\label{vii:prop:naturalidad-proyeccion-absorcion}
For a source isometry \(R=R_m^J\), one has
\(P_{{\rm abs},m+1}R=RP_{{\rm abs},m}\) exactly when
\[
 R(\ker\mathcal A_{\partial,m})
       \subset\ker\mathcal A_{\partial,m+1},\qquad
 R((\ker\mathcal A_{\partial,m})^\perp)
       \subset(\ker\mathcal A_{\partial,m+1})^\perp.
\]
Under these conditions at every level, the projections define an orthogonal
projection on the completion of the isometric inductive limit of the source
spaces.
\end{proposition}
\begin{proof}
Each source is decomposed into its components in the kernel and its
orthogonal complement. The two inclusions imply that applying the
projection before or after refinement respectively preserves the first
component and annihilates the second. Conversely, the operator equality
applied to each component gives the two inclusions.
In the inductive limit, the equality allows an operator to be defined on
the union of the levels independently of the representative. All projections
have norm at most one, so the operator extends to the completion.
Its self-adjointness and idempotence identities, valid on the dense union,
pass by continuity to the completion.
\end{proof}

To prolong a Green response to a completion as well, its domains and the
corresponding continuity estimates are retained. Chronological propagation
with finite support and factorization of the response through the boundary
image, proved in \ref{vii:thm:observador-factorizacion}, thus retain their
respective domains throughout refinement.

The physical realization of emission and absorption composes these operators
with the matter currents, action, and causal conditions of the sector.
The intertwining must preserve the route involution, source--field pairing,
transported action, and boundary conditions at each level.
The preceding construction provides group reversal, the Green operators of the
chronological channel, temporal decomposition, and an explicit absorbing
projection; spatial realizations use their own maps and causality conditions.
